\chapter{2011 Nobel Prize in Economics}
\textbf{Laureates: }Christopher Sims (USA), Robert Lucas Jr. (USA)

\section*{Reason for Award}
For their empirical research on cause and effect in the macroeconomy

\section{What They Did}
Robert Lucas revolutionized macroeconomics by applying the rational expectations
hypothesis, showing that systematic monetary policy cannot systematically affect
real variables like output and employment. Born in Gary, Indiana in 1937, Lucas
earned his doctorate at Chicago and spent most of his career at the University
of Chicago, where he became a leading figure in the new classical macroeconomics
movement.

His critique demonstrated that econometric models ignoring expectations formation
were unreliable for policy analysis. Lucas argued that agents form expectations
rationally, using all available information including policy rules. This implied
that only unexpected policy changes affect real variables. The Lucas critique
showed that estimated relationships in econometric models would change when
policy changes, because agents adjust their behavior in response to new policy
rules.

Christopher Sims developed vector autoregression methodology for analyzing cause
and effect in macroeconomics. Born in Salem, Oregon in 1942, Sims earned his
doctorate at MIT and has been a professor at Princeton.

VAR models treat all variables as endogenous, allowing researchers to identify
structural shocks without imposing strong theoretical restrictions. Sims' impulse
response analysis showed how economic variables respond to unexpected changes in
monetary policy, fiscal policy, and other shocks. This approach became known as
"identification without theory."

Together, their work established modern macroeconometrics and the study of
business cycle dynamics through empirical identification of causal relationships.
Lucas provided the theoretical foundation while Sims provided the empirical
tools for testing macroeconomic theories.

Lucas's work on the Lucas critique showed that econometric relationships change
when policy changes, because agents adjust their expectations. This undermined
the Keynesian models that dominated policy analysis before the 1970s.

Sims' VAR methodology provided a flexible alternative to structural models,
allowing economists to identify causal effects from data without strong
theoretical priors. His work showed how to interpret VAR estimates in terms
of structural shocks using minimal assumptions.

\section{Impact on Economic Thinking}
Lucas and Sims transformed macroeconomic theory and empirical practice. Lucas'
rational expectations revolution challenged Keynesian economics, showing that
policy effectiveness depends on expectations formation. This led to the
development of new classical macroeconomics and real business cycle theory.

Central banks incorporated Lucas' insights into their policy frameworks,
recognizing that credibility and expectations management matter. The idea
that policy effects depend on how agents expect policy to evolve became
central to modern monetary policy design.

Sims' VAR methodology provided empirical tools for testing macroeconomic theories
and analyzing policy effects. Impulse response functions became standard for
communicating how economies respond to shocks. Their work established the
importance of identifying structural shocks rather than merely correlating
variables.

This methodological advance improved empirical macroeconomics across many
subfields. The Lucas critique remains a fundamental principle in policy analysis,
reminding economists that model parameters depend on policy regimes. Sims'
approach to identification became standard in applied macroeconomics.

The rational expectations revolution changed how economists think about policy.
It showed that systematic policy cannot systematically affect real variables,
shifting focus to the role of expectations and credibility in macroeconomic
management.

\section{Modern Perspectives Today}
Sims and Lucas' contributions remain foundational in macroeconomic research and
policy. VAR models are widely used for forecasting, policy analysis, and
structural interpretation. Modern macroeconomics incorporates rational
expectations as a standard assumption.

The methodology developed by Sims has been extended to high-frequency
identification, narrative approaches, and nowcasting. Central banks use these
techniques to evaluate policy rules and assess the effects of unconventional
monetary policies like quantitative easing.

The 2008 financial crisis renewed interest in macroeconomic shocks and
propagation mechanisms, areas where both scholars made key contributions.
Research on forward guidance, quantitative easing, and fiscal multipliers
builds on their work.

The rational expectations hypothesis has been extended to incorporate bounded
rationality and learning, addressing criticisms about perfect information
assumptions. Modern macroeconometric practice continues to rely on the
identification strategies and empirical tools pioneered by Sims and Lucas.

Their work established the standards for empirical macroeconomics that continue
to guide research today. The emphasis on credible identification and the
importance of expectations remains central to macroeconomic analysis.
